% 第 46 章。数字与教学模型：checks/check_ch46.py；图：figsrc/fig_ch46.py。
% 史料：floating_facts（取数脚本 data/fetch/floating46.py）；汇率与物价：floating_bis_1959_1985（同一脚本）。
\chapter{布雷顿森林解体与浮动汇率}\label{ch:floating}

\begin{chapterguide}
\guidequestion{美元与黄金为什么脱钩？浮动汇率为什么会“超调”？}
\begin{guidegoals}
\item 追踪 1971 年 8 月美元停止兑换黄金到 1973 年主要货币转向浮动的过程；
\item 推导多恩布什（1976）的汇率超调模型；
\item 说明浮动以后汇率波动大于基本面的现象（Meese \& Rogoff 1983）；
\item 描述欧洲的“蛇形浮动”与早期汇率合作。
\end{guidegoals}
% 据: nx_date, b73_joint, db_journal, mr_journal, b73i_snake
\guideprereq{\chref{ch:brettonwoods}\enspace\nameref{ch:brettonwoods}；\chref{ch:mundell}\enspace\nameref{ch:mundell}}
\end{chapterguide}

1971 年 8 月 15 日是星期日。晚上 9 点，尼克松（Richard Nixon，1913—1994）在白宫的椭圆形办公室发表讲话，广播与电视同时直播\cite{nixon1971,loc46}。他说，新经济政策要在三条战线上作战，创造更多更好的工作，制止生活费用上涨，保护美元不受国际货币投机者的攻击\cite{nixon1971}。他宣布冻结全国的物价与工资 90 天，削减 10\% 的对外经济援助\cite{nixon1971}。
% 据: nx_date, nx_sunday, nx_note, nx_fronts, nx_freeze, nx_aid, loc_nixon

美元的事放在后面。他说，过去 7 年里平均每年有一次国际货币危机，得益的是国际货币投机者\cite{nixon1971}。他已指示财政部长康纳利（John B. Connally，1917—1993）暂时停止美元兑换黄金或其他储备资产\cite{nixon1971,loc46}。他还宣布对进口商品加征 10\% 的附加税，说这是为了不让美国产品因为“不公平的汇率”吃亏\cite{nixon1971}。
% 据: nx_seven, nx_gainers, nx_suspend, nx_surcharge, nx_unfair, loc_connally

这一步意味着什么，他也替听众作了回答。他说，要打消人们对所谓贬值的担心\cite{nixon1971}。买外国汽车或出国旅行，美元也许能买的东西少一点\cite{nixon1971}。在美国买美国货的绝大多数美国人，手里的美元“明天和今天一样值钱”\cite{nixon1971}。
% 据: nx_bugaboo, nx_abroad, nx_tomorrow

一年半以后，主要货币不再盯住美元。1973 年 3 月 19 日外汇市场重新开放，欧洲共同体的 6 个国家对美元共同浮动，英国与意大利各自浮动，加拿大、日本与瑞士也继续浮动\cite{bis1973}。黄金与美元的联系从此断了，汇率天天在变。美元与黄金为什么脱钩？浮动汇率为什么会“超调”？
% 据: b73_reopen, b73_joint, b73_uk, b73_others

\section{当时的世界}\label{sec:floating-world}

\chref{ch:brettonwoods}讲过这个体系怎样运转。1946 年起，金价是每盎司 35 美元，盎司是称黄金用的重量单位。美国是唯一按黄金钉住本国货币、买卖黄金的国家，其他国家都对美元定下平价，并在平价上下 1\% 以内干预汇率。各国货币经由美元间接地连着黄金（本书的概括）。所谓黄金窗口，就是外国央行拿美元向美国换黄金的那扇门（本书的解读）。
% 据: b72i_margins, b72_smithsonian

危机有它的来路。国际清算银行（BIS）在 1972 年的年报里回顾说，1970 年美国的短期资本净流出 76 亿美元，官方结算差额出现了史无前例的 98 亿美元赤字\cite{bis1972}。美国的货币政策在 1969—1970 年由紧转松，欧洲的政策却朝相反的方向转\cite{bis1972}。差别最大的是联邦德国。1970 年 3 月，美联储正在放松，联邦银行却提高了贴现率\cite{bis1972}。
% 据: b72_deficit, b72_settlements, b72_swing, b72_germany

从美国流出的钱，主要流进了联邦德国。1970 年的后 9 个月，联邦德国的储备增加了 58 亿美元\cite{bis1972}。BIS 说，到 1970 年底，美国的基本收支已经很弱，扶持复苏的货币政策又必然带来更多的资金外流，局面一触即发\cite{bis1972}。这时各国的汇率已经很难守住。BIS 说，1969 年 8 月到 1970 年 6 月，法国法郎贬了值，马克先浮动后升值，加元也浮动了\cite{bis1972}。马克是联邦德国的货币，\chref{ch:brettonwoods}讲过它 1969 年的浮动。1971 年 5 月，马克又一次浮动，荷兰盾也一起浮动\cite{bis1972,bis1973}。
% 据: b72_reserves70, b72_weak, b72_outflows, b72_explosive, b72_franc, b72_mark, b73_floats

8 月又起了风浪。BIS 说，8 月的危机始于 6 日，美国国会的一个小组委员会在报告里要求各国货币对美元普遍调整汇率\cite{bis1972}。报告还说，不然美国也许只能单方面停止美元的兑换\cite{bis1972}。BIS 评论说，在别国听来，这像是从内部直接透出的消息\cite{bis1972}。此后一周，逃离美元的资金使其他国家的储备增加 37 亿美元，流进瑞士法郎的最多，有 17 亿美元\cite{bis1972}。
% 据: b72_august, b72_report, b72_unilateral, b72_horse, b72_flight

美国自己的储备也在流失。BIS 说，1971 年的官方结算赤字主要靠对外国官方机构的负债增加 274 亿美元来弥补\cite{bis1972}。导致美元正式停止兑换的，却是储备的流失\cite{bis1972}。它同堆积起来的债务相比不算大，同剩下可以流失的储备相比却很大，而且一个季度比一个季度严重\cite{bis1972}。
% 据: b72_liabilities, b72_losses, b72_quarter

谁作的决定？8 月 2 日，尼克松几次会见康纳利与管理和预算局局长舒尔茨（George P. Shultz，1920—2021）\cite{frus46a,loc46}。《美国对外关系文件》（FRUS）的编者据录音说，康纳利是这套方案的主要设计者，他那天几乎讲全了 8 月 15 日要宣布的要点\cite{frus46a}。在国际方面，他主张停止兑换、让美元浮动、加征 10\% 的进口附加税\cite{frus46a}。他认为兑换维持不到 1972 年的大选，多半年底以前就得停止\cite{frus46a}。
% 据: f164_aug2, f164_architect, f164_elements, f164_intl, f164_election, loc_shultz

康纳利把国内与国际两方面连在一起。他说，国内的措施表明美国愿意为保卫美元采取艰难的步骤\cite{frus46a}。附加税他甚至一度说成国内方案的一部分，每年能给预算带来 40 亿美元\cite{frus46a}。尼克松倾向于必要时 8 月就行动\cite{frus46a}。他交代，需要和外国商量，就找基辛格（Henry Kissinger，1923—2023）的渠道\cite{frus46a,loc46}。编者转述，他还说在任何情况下都不要与国务院商量，因为国务院代表的是外国政府\cite{frus46a}。
% 据: f164_linked, f164_budget, f164_leaning, f164_kissinger, f164_state, loc_kissinger

8 月 11 日，尼克松对舒尔茨说，联邦储备委员会主席伯恩斯（Arthur F. Burns，1904—1987）对此“非常紧张”\cite{frus46a,loc46}。舒尔茨提醒他，只关闭黄金窗口，别国若干预市场维持本币，美国未必得到想要的汇率变化\cite{frus46a}。他主张立即关闭黄金窗口，加征临时的进口税，也就是变相贬值，然后谈判\cite{frus46a}。
% 据: f165_nervous, f221_burns, f165_window, f165_devaluation, loc_burns

日子是被市场催出来的。8 月 12 日中午，康纳利在电话里对尼克松说，他们“正在失去主动”\cite{frus46a}。当天下午，舒尔茨汇报了 8 月 2 日以来各国拿美元兑换的规模\cite{frus46a}。尼克松说，主要的人都去戴维营，下星期一就能准备好行动\cite{frus46a}。他仍觉得全部方案最好 9 月 7 日宣布，黄金窗口的关闭可以先由康纳利低调宣布\cite{frus46a}。
% 据: f165_aug12, f165_initiative, f165_afternoon, f165_conversions, f165_monday, f165_sept7, f165_lowkey

他也估计到了阻力。尼克松认为，多数顾问会赞成，伯恩斯与财政部次长沃尔克（Paul A. Volcker，1927—2019）却会对方案的国际部分有异议，得把他们拉过来\cite{frus46a,loc46}。他让舒尔茨去说服伯恩斯，要伯恩斯明白只有两三种选择\cite{frus46a}。
% 据: f165_advisers, f165_international, f165_burns, f165_burnsname, loc_volcker

8 月 13 日下午，尼克松在戴维营与经济顾问及白宫幕僚开会\cite{frus46a}。按总统日志，到会的有 10 人\cite{frus46a}。其中有康纳利、伯恩斯、舒尔茨、沃尔克与斯坦（Herbert Stein，1916—1999）\cite{frus46a,loc46}。会议断断续续开到 15 日上午 11 点半\cite{frus46a}。国务院与国家安全委员会没有派人，开会时也不知情\cite{frus46a}。基辛格后来在回忆录里说，这个对外交极为重要的决定，国务卿与国家安全顾问事先都没有被征询\cite{frus46a}。
% 据: f168_date, f168_time, f168_diary, f168_until, f168_absent, f168_unaware, f168_kissinger, f221_volcker, loc_stein

会上反对的声音很弱。编者说，霍尔德曼的会议笔记记下了伯恩斯想让黄金窗口开着的努力，孤单而没有成功\cite{frus46a}。斯坦 1998 年接受编者访谈时说，与会者并不清楚，关闭黄金窗口是不是意味着美元浮动\cite{frus46a}。他还说，国务院与国家安全委员会没有人到会，会上很少考虑国际政治的后果\cite{frus46a}。到 8 月 14 日，尼克松才说要通知国务卿，问他是否同意削减 10\% 的对外援助\cite{frus46a}。对盟国影响最大的决定，管外交的部门却没有参与（本书的解读）。
% 据: f168_burns, f168_stein, f168_politics, f168_rogers

代价落在谁身上？附加税落在向美国出口的国家身上。尼克松说，这个税不针对任何国家，不公平的待遇一结束，它也就结束\cite{nixon1971}。BIS 记下了取消的条件。别国货币要充分升值，阻碍美国出口的壁垒要放松，防务负担要分担得更好，国际货币制度的改革要达成协议\cite{bis1972}。
% 据: nx_notdirected, nx_end, b72_conditions

浮动也有先例。\chref{ch:mundell}讲过加拿大 1950 年代的浮动。1970 年 5 月 31 日，加拿大的财政部长宣布外汇基金暂时停止买入美元，加元再次浮动\cite{powell2005}。按 Powell 的叙述，政府作这个决定并不情愿\cite{powell2005}。
% 据: pw_may31, pw_reluctant

\section{从史密森协定到全面浮动}\label{sec:floating-smithsonian}

8 月 23 日是星期一，关闭了一周的欧洲主要外汇市场重新开放\cite{bis1972}。除法国以外，欧洲国家都让本币浮动\cite{bis1972}。法国只对贸易与官方交易维持原来的汇率，另设一个金融外汇市场\cite{bis1972}。日本的市场没有关闭，日本银行在干预上限继续买入美元，买了两个星期\cite{bis1972}。
% 据: b72_reopen, b72_float, b72_france, b72_japan

浮动并不自由。BIS 说，各国一面管制资本流入，一面照旧在市场上买入美元\cite{bis1972}。8 月底到 12 月底，联邦德国与日本的总储备分别增加 17 亿与 28 亿美元\cite{bis1972}。BIS 列举了不愿完全自由浮动的理由。一是出口商等利益集团在政治上强烈反对本币升值，经济放缓又加强了这种反对\cite{bis1972}。二是各国对美元升值的幅度不一，竞争力的变化不只是对美元\cite{bis1972}。
% 据: b72_controls, b72_reserves, b72_exporters, b72_slowed, b72_unequal

美国起初并不急。BIS 说，其他主要国家都希望尽快谈出一套新的固定汇率，美国一度不愿参加\cite{bis1972}。浮动汇率对企业信心的不利影响一显露，这个态度很快就变了\cite{bis1972}。
% 据: b72_reluctant, b72_confidence

升多少，各方算的不一样。美国认为它的国际收支需要改善 130 亿美元，国际货币基金组织与经济合作与发展组织算出来是 90 亿美元左右\cite{bis1972}。11 月 30 日，康纳利从罗马给白宫发电报，附上美国向十国集团提出的方案\cite{frus46a}。方案要求以 1971 年 5 月 1 日为基点，其他经合组织国家的货币对美元加权平均至少升值 11\%，黄金的美元价格不变\cite{frus46a}。
% 据: b72_13, b72_9, f211_date, f211_11, f211_gold

电报的理由是美国的赤字太大。它说，1971 年头 9 个月，美国的基本赤字按年率估计在 100 亿美元以上\cite{frus46a}。过去 7 个月，连贸易也出现了逆差，按年率约 30 亿美元\cite{frus46a}。方案还说，美国不能指望恢复美元兑换储备资产，也不能提供汇率担保\cite{frus46a}。作为交换，美国取消附加税，条件是贸易上的让步与防务开支的分担\cite{frus46a}。
% 据: f211_deficit, f211_tradedef, f211_noconv, f211_surcharge, f211_trade, f211_exporters

罗马没有达成协议，美国在会上同意考虑一个改善不到 130 亿美元的调整，也同意考虑美元贬值\cite{bis1972}。12 月 17—18 日，十国集团的部长在华盛顿的史密森学会达成协议\cite{frus46a}。美元对黄金贬值 7.9\%，日元对美元升值 16.9\%，马克升值 13.6\%（\tabref{tab:floating-smithsonian}）\cite{bis1972}。对美元的波动幅度从平价上下各 1\% 放宽到各 2.25\%，上下限之间的距离达到 4.5\%\cite{bis1972}。
% 据: b72_rome, b72_noagree, f221_date, b72_smithsonian, b72t_usd, b72t_jpy, b72t_dem, b72i_margins

美元的平价要等国会通过才能改，别国也就没有马上改黄金平价\cite{bis1972}。要改平价的国家先对美元宣布与新平价相当的“中心汇率”，不改黄金平价的国家则对美元定“中间汇率”\cite{bis1972}。幅度放宽以后，央行不必等汇率离开中心 1\% 就入市，可以让它走得更远（本书的解读）。
% 据: b72_congress, b72_central, b72_middle

新汇率下，日元很快站到了升值的一边。按表中日元的中心汇率 308 日元，上下各 2.25\%，下限约是 301.1 日元\cite{bis1972}。1972 年 7 月到 1973 年 1 月，日元的月平均都在 301 日元上下，贴着这条线（本书的计算）\cite{bisdata46}。市场押的是日元还要升值（本书的解读）。
% 据: b72t_jpy, b72i_margins, bisdata_jpy, bisdata_jpy7301

\begin{table}[tbp]\centering\small
\caption{1971 年 12 月 18 日的汇率调整（\%）与新汇率（每美元合多少本币）}\label{tab:floating-smithsonian}
\begin{tabularx}{\linewidth}{@{}lYYY@{}}
\toprule
货币 & 对黄金 & 对美元 & 新的中心汇率或中间汇率\\
\midrule
日元 & $+7.7$ & $+16.9$ & 308.00\\
瑞士法郎 & $+7.1$ & $+13.9$ & 3.84\\
奥地利先令 & $+5.1$ & $+11.6$ & 23.30\\
德国马克 & $+4.6$ & $+13.6$ & 3.2225\\
比利时法郎 & $+2.8$ & $+11.6$ & 44.8159\\
荷兰盾 & $+2.8$ & $+11.6$ & 3.2447\\
英镑 & — & $+8.6$ & 2.60571\\
法国法郎 & — & $+8.6$ & 5.1157\\
意大利里拉 & $-1.0$ & $+7.5$ & 581.50\\
瑞典克朗 & $-1.0$ & $+7.5$ & 4.8129\\
美元 & $-7.9$ & — & —\\
\bottomrule
\end{tabularx}
\datasource{国际清算银行 1972 年年报第 30 页的表 Realignment of leading currencies，1971 年 5 月 5 日到 12 月 18 日的变化。英镑一行的新汇率是每英镑合多少美元。英镑与法国法郎对黄金的平价没有变。先令是奥地利的货币，克朗是瑞典的货币。中文是本书的翻译，不是 BIS 的官方译文。}
% 据: b72t_title, b72t_head, b72t_jpy, b72t_chf, b72t_ats, b72t_dem, b72t_bef, b72t_nlg, b72t_gbp, b72t_frf, b72t_itl, b72t_sek, b72t_usd, b72t_note
\end{table}

康纳利在会后的记者会上说，美元对其他经合组织国家的货币平均贬值 12\%，政府将请国会授权把黄金价格提高到每盎司 38 美元\cite{frus46a}。下午 5 点 48 分，尼克松来到史密森学会，称这是“历史上最重要的货币协定”\cite{nixon1971b}。他说，有人会问谁赢谁输，答案是整个自由世界都赢了\cite{nixon1971b}。
% 据: f221_12, f221_38, nxs_note, nxs_most, nxs_won

协定要靠各国自己落实。法国 12 月 21 日宣布了经常账户交易的新平价，金融法郎却不受新幅度的约束\cite{frus46a}。1972 年 4 月 3 日，尼克松签署了改变美元平价的法案，美国 5 月 5 日正式通知国际货币基金组织\cite{bis1972}。
% 据: f221_france, b72_bill, b72_imf

新汇率并不稳。1972 年 6 月 23 日起，英国停止按固定的价格买卖美元，英镑开始浮动\cite{bis1973}。BIS 1973 年回顾说，汇率的调整一再拖延，总要等到市场危机的不可抗力才启动\cite{bis1973}。1973 年 1 月 22 日，瑞士法郎升到对美元的上限，瑞士当局连买两天美元以后让法郎浮动，以免违背反通胀的目标\cite{bis1973}。
% 据: b73_sterling, b73_delayed, b73i_jan22, b73i_swiss

2 月的头 7 个交易日，十国集团的货币当局买入约 80 亿美元，其中联邦德国近 60 亿，日本占其余的一半\cite{bis1973}。到 2 月 9 日星期五，局面已经撑不住，下一周开始时多数外汇市场关闭\cite{bis1973}。美元第二次贬值 10\%，日元与意大利的商业里拉随即浮动\cite{bis1973}。2 月 14 日市场重开以后，日元在美元贬值之外又很快升值 5\%\cite{bis1973}。
% 据: b73i_feb, b73i_untenable, b73i_devalue, b73i_yen, b73i_yen5

宣布的日子，FRUS 记作 2 月 12 日晚上，日本银行的行史也记作 12 日，BIS 记作 13 日，差的大概是时差（本书的解读）\cite{frus46b,boj1986,bis1973}。宣布贬值的是已任财政部长的舒尔茨\cite{frus46b}。他说，美国坚定地期望，日元浮动到一个不靠政府大量干预的位置\cite{frus46b}。尼克松在一份文件上批注，提醒要强调 1971 年 8 月 15 日讲过的那一点，贬值不减少美元在国内的价值\cite{frus46b}。
% 据: f12_evening, boj_feb12, b73i_devalue, f12_yen, f12_note

3 月 2 日，尼克松在记者会上说，“不会再有一次贬值”\cite{frus46b}。前一天，联邦银行一天就买进 27 亿美元，外汇市场又关闭了\cite{bis1973}。联邦德国 2 月头 7 个交易日买入近 60 亿美元，3 月 1 日又买 27 亿美元，马克随后升值，这些美元按马克计都缩了水（本书的解读）。欧洲共同体的 6 国随后商定，彼此的汇率维持在名义交叉汇率上下 2.25\% 以内，对美元则听由市场决定，马克另外升值 3\%\cite{bis1973}。
% 据: f14_date, f14_nodeval, b73i_march1, b73_joint, b73_margin, b73_revalue

日子又有出入。BIS 的年报一处说 3 月 12 日作出决定，另一处说 11 日达成安排，FRUS 与日本银行的行史都记作 11 日\cite{bis1973,frus46b,boj1986}。3 月 16 日，十国集团的部长与丹麦、爱尔兰、卢森堡、瑞士及几个国际组织的代表在巴黎开会，认可了共同体 3 月 11 日的安排\cite{frus46b}。会上同意，必要时可以由中央银行干预来管理浮动，美国也可能参与\cite{bis1973}。BIS 说，这一次人们认为回到美元与其他主要货币之间的固定汇率已是徒劳\cite{bis1973}。3 月 15 日晚上，西蒙在电话里对基辛格说，6 个国家共同浮动，英国、意大利、加拿大等都在各自浮动\cite{frus46b}。
% 据: b73_joint, b73i_concluded, f35_march11, boj_march11, f35_paris, b73_g10, b73i_futile, f35_simon

浮动并没有带来平静。BIS 1974 年说，到那时为止浮动的主要特点，是汇率宽幅而不规则的变化\cite{bis1974}。美元对欧洲货币先跌得远远超过 10\% 的贬值，又连涨约六个月，不但收复了跌幅，还抹掉了大部分贬值\cite{bis1974}。此后美元又转弱\cite{bis1974}。
% 据: b74_erratic, b74_dollar, b74_weak

BIS 的判断是，市场汇率波动这样大，不得不大规模地官方干预，浮动的经验难说令人满意\cite{bis1974}。所以国际货币基金组织的二十国委员会继续寻求“稳定但可调整的平价”，这项工作却没有成功\cite{bis1974}。它还说，通胀加剧了汇率的不稳，浮动汇率也是通胀加速的一个因素\cite{bis1974}。货币贬值的国家，进口价格上涨直接推高物价\cite{bis1974}。汇率的大起大落又削弱了人们对货币的信心，使人们偏爱各种实物资产\cite{bis1974}。
% 据: b74_intervention, b74_c20, b74_inflation, b74_imports, b74_confidence

同一年，油价也涨了起来。BIS 1974 年把油价的猛涨列为过去一年最引人注目的事\cite{bis1974}。它的结论是，头一年浮动的经验表明，总的来说浮动让货币稳定更难实现，也更难维持\cite{bis1974}。油价在\chref{ch:oil}。
% 据: b74_oil, b74_complicated

浮动从权宜之计变成了制度。1976 年 1 月，国际货币基金组织理事会的临时委员会在牙买加完成协议，修改后的协定第四条让各国自选汇率安排，黄金的官方价格也取消了\cite{bis1976}。按 BIS 的说法，华盛顿与牙买加的协议正式结束了布雷顿森林的世界，却没有产生一套真正有约束力的新规则\cite{bis1976}。
% 据: b76_jamaica, b76_article4, b76_gold, b76_end

日本与欧洲走了两条路。日本守着 360 日元，守到守不住，下面的专栏讲这段经过与它的代价（本书的概括）。共同体国家则想在彼此之间维持固定汇率，对美元一起浮动，这就是第二个专栏里的蛇形浮动（本书的概括）。

\begin{case}{日本 · 1971—1973 · 360 日元的终结}
尼克松的讲话传到东京，是日本时间 8 月 16 日（星期一）上午 10 点，东京外汇市场已经开市\cite{boj1986}。日本银行总裁佐佐木直（生卒年未详）与大藏大臣水田三喜男（1905—1976）随即发表谈话，坚持现行平价\cite{boj1986,loc46}。日本银行指出关闭市场的问题，表示关闭不妥\cite{boj1986}。最后由大藏大臣裁断，17 日以后照常开市\cite{boj1986}。
% 据: boj_tokyo, boj_keep, boj_bank, boj_decision, loc_mizuta

代价落在谁身上？首先是外汇储备。8 月 16 日到 27 日的 11 天里（星期日除外），平衡用的美元买入近 40 亿美元，约相当于当时 4 个月的进口，外汇储备从不到 80 亿美元猛增到 120 亿美元\cite{boj1986}。卖出美元的主要是提前收汇的出口商与调整头寸的银行，非居民的投机资金比较少\cite{boj1986}。8 月 28 日（星期六）日元转向浮动，当天升到 341.30 日元\cite{boj1986}。
% 据: boj_11days, boj_sellers, boj_exporters, boj_float, boj_341

其次是出口产地。按经济企画厅 1972 年白皮书的统计，输出比率超过 10\% 的中小企业产地全国有 120 处，其中对美国的出口占 50.6\%\cite{epa1972}。为保住出口数量，这些产地压低以日元计的出口价格，再靠压低计件工钱、延长劳动时间来消化\cite{epa1972}。经济企画厅指出，产地的工人多是家庭主妇与农民，容易转成减薪、加班、转业或歇业\cite{epa1972}。
% 据: epa_120, epa_cuts, epa_workers

日本银行后来在行史里检讨这段经过。它说，花 11 天大量买进美元，后来招致不少批评\cite{boj1986}。最大的问题，是对日元升值的影响作了极度悲观的判断\cite{boj1986}。按行史的转述，佐佐木后来也承认，当时把升值的影响看得过大，对国际货币问题的解决想得有些简单\cite{boj1986}。
% 据: boj_criticism, boj_pessimism, boj_sasaki

1973 年 2 月 10 日（星期六），大藏大臣爱知揆一（1907—1973）裁断，日本先于他国关闭了东京外汇市场\cite{boj1986,loc46}。日本时间 2 月 13 日各国达成协议，日本因为预算案正在国会审议、难以改变平价，转向浮动\cite{boj1986}。爱知的谈话说，政府希望尽早回到固定汇率\cite{boj1986}。日元随即升到 263 日元台\cite{boj1986}。
% 据: boj_feb10, boj_feb13, boj_aichi, boj_263, loc_aichi
\end{case}

\begin{case}{欧洲 · 1972—1976 · 蛇形浮动}
欧洲共同体有自己的计划。1972 年 3 月 21 日理事会的决议要求，作为走向共同体货币区的第一步，各国央行逐步缩小成员国货币之间的差距，到 1972 年 7 月 1 日不超过 2.25\%（本书据法文本翻译）\cite{ojc1972}。4 月 24 日起，计划付诸实施，后来常被称为“隧道里的蛇”\cite{bis1973}。
% 据: oj_date, oj_first, oj_225, b73i_snake

规则是这样的。任何两种成员国货币，一段时间里的波动与同一时点的差距，都不得超过成员国货币与美元之间可能的幅度，即 4.5\% 与 2.25\%\cite{bis1973}。美元给的幅度是隧道，成员国之间的窄幅是隧道里的蛇（本书的解读）。英国与丹麦 5 月 1 日加入，挪威 5 月 23 日加入\cite{bis1973}。
% 据: b73i_rule, b73i_join

蛇很难养。英镑 6 月离开，丹麦同时退出，10 月 11 日又回来\cite{bis1973}。意大利先得到临时豁免，后来也退出了\cite{bis1973}。1973 年 3 月以后，蛇对美元共同浮动，没有了隧道（本书的解读）。欧洲货币合作基金 1973 年 4 月成立，6 月 1 日开始运作\cite{bis1973}。
% 据: b73i_leave, b73_denmark, b73i_italy, b73_emcf

法国进进出出。它 1974 年 1 月 21 日退出共同浮动\cite{bis1974}。1975 年 7 月按原来的中心汇率回来，1976 年 3 月 15 日在储备累计损失 30 亿美元以后再次退出\cite{bis1976}。瑞士不在蛇里。按 Rich 的叙述，1973 年 1 月新的投机攻击逼近，瑞士当局不情愿地让法郎浮动\cite{rich2003}。
% 据: b74_france, b76_return, b76_leave, rich_float73

1978 年法郎升得太猛，瑞士国家银行在 9 月考虑过一个目标区，马克的价格在每 100 马克 84 到 100 瑞士法郎之间\cite{rich2003}。它没有采用目标区，10 月初改为临时的汇率目标（\chref{ch:monetarism}）\cite{rich2003}。按 Rich 的说法，政策转向以后，法郎开始回落到比较现实的水平\cite{rich2003}。瑞士国家银行在 1979 年的年报里说，短暂的汇率目标试验以后，它又回到让基础货币平稳增长的策略\cite{rich2003}。
% 据: rich_zone, rich_zone84, rich_80, rich_depreciate, rich_return
\end{case}

\section{超调模型}\label{sec:floating-model}

浮动以后，汇率比物价动得快得多（\secref{sec:floating-test}）。多恩布什（Rudiger Dornbusch，1942—2002）1976 年的论文给了一个解释\cite{dornbusch1976,loc46}。按他的摘要，汇率最初的\textbf{超调}（overshooting）来自不同市场调整速度的差别\cite{dornbusch1976}。资产市场随时出清，商品的价格却慢慢调整。
% 据: db_abstract, db_journal, loc_dornbusch

先用文字说一遍。货币突然增加，物价一时不动，实际货币多了，本国利率就要下降。资本完全流动，利率低于外国，持有本币的人就要得到补偿，只能是预期本币以后升值。本币要能升值，现在就得先贬过头，落到比长期水平更低的地方（本书的解读）。汇率先跳过头，再慢慢回来，这就是超调。
% 据: db_abstract

下面沿用\chref{ch:external}的记号，只是把 $e$、$p$、$m$ 都取成对数，分别是本币表示的外币价格、本国价格水平与货币量。外国的价格水平固定，实际汇率的对数就是 $q=e-p$，即式~\eqref{eq:external-real}取对数。本国是小国，产出固定在充分就业的水平。模型有三个假设\cite{dornbusch1976}。
% 据: db_small, db_fixed, db_foreign

第一，资本完全流动，式~\eqref{eq:external-uip}取对数就是 $i=i^{f}+x$，$x$ 是预期的本币贬值率。第二，汇率高于长期水平 $\bar e$，人们就预期它回落，$x=\theta(\bar e-e)$。第三，货币市场出清，$m-p=k-\lambda i$，$k$ 包含固定的产出\cite{dornbusch1976}。多恩布什用 $r$ 与 $r^{*}$ 表示本国与世界的利率，本书的 $r^{*}$ 留给中性利率，所以改写成 $i$ 与 $i^{f}$，他的式（1）至（3）与这里的三个式子相同（本书的核对）。
% 据: db_eq1, db_eq2, db_eq3

长期里 $e=\bar e$，$i=i^{f}$，物价是 $\bar p=m-k+\lambda i^{f}$，即多恩布什的式（5）\cite{dornbusch1976}。货币长期中性，购买力平价成立，货币增加 $\mathrm{d}m$，$\bar p$ 与 $\bar e$ 也增加 $\mathrm{d}m$（\chref{ch:external}、\chref{ch:monetarism}）。短期里物价不动，把三个假设合起来得到
\begin{keyeq}
e=\bar e-\frac{1}{\lambda\theta}\,(p-\bar p),\qquad
\frac{\mathrm{d}e}{\mathrm{d}m}=1+\frac{1}{\lambda\theta} .\label{eq:floating-overshoot}
\end{keyeq}
左式是多恩布什的式（6），右式是他的式（16）\cite{dornbusch1976}。货币增加的那一刻，$\bar e$ 上升 $\mathrm{d}m$，$p$ 不变，$e$ 却上升 $(1+1/\lambda\theta)\,\mathrm{d}m$，跳过了长期水平（本书的推导）。利率降了 $\mathrm{d}m/\lambda$，本币必须预期升值，才能抵消这个利差（本书的推导）。
% 据: db_eq5, db_eq6, db_eq16, db_eq16i

超调的大小取决于两个参数。多恩布什说，货币需求对利率越敏感，超调越小，因为同样的货币增加只让利率下降一点\cite{dornbusch1976}。预期系数 $\theta$ 越大，汇率只要稍微高出长期水平，就足以产生所需的升值预期，超调也越小（本书的解读）。
% 据: db_dampen

物价若能随时调整，就不会超调。货币增加 $\mathrm{d}m$，物价当时就上涨 $\mathrm{d}m$，实际货币不变，利率也不变，汇率只要一次贬值 $\mathrm{d}m$ 就到位了（本书的推导）。这是弹性价格的货币模型。超调全靠物价跟不上，多恩布什的标题里那个“动态”就在这里（本书的解读）。
% 据: db_title

以后物价慢慢上涨。多恩布什设物价按商品的超额需求上涨，速度是 $\pi$\cite{dornbusch1976}。需求随本国商品变便宜（$e-p$ 上升）而增加，系数是 $\delta$，随利率上升而减少，系数是 $\sigma$\cite{dornbusch1976}。预期要与模型自己的路径一致，$\theta$ 必须等于经济实际收敛的速度，于是\cite{dornbusch1976}
\begin{equation}
\theta=\frac{\pi}{2}\Bigl(\frac{\sigma}{\lambda}+\delta\Bigr)+\sqrt{\frac{\pi^{2}}{4}\Bigl(\frac{\sigma}{\lambda}+\delta\Bigr)^{2}+\frac{\pi\delta}{\lambda}} .\label{eq:floating-theta}
\end{equation}
此后 $e-\bar e$ 与 $p-\bar p$ 都按 $\mathrm{e}^{-\theta t}$ 缩小，汇率一路回落，物价一路上涨，利率回到 $i^{f}$（本书的推导）。多恩布什说，货币需求对利率越不敏感，商品需求对利率与相对价格越敏感，经济收敛得越快\cite{dornbusch1976}。
% 据: db_eq8, db_eq7, db_eq15, db_faster

这个模型与蒙代尔—弗莱明模型问的不是一件事。那里物价固定，问的是货币政策在浮动汇率下能不能改变产出。多恩布什把产出固定，让物价慢慢调整，问的是汇率怎样从短期走到长期（本书的概括）。

实际汇率在这里起了作用。货币增加的那一刻，$e$ 跳上去，$p$ 不动，$q=e-p$ 也跳上去，本国商品变便宜了。需求增加，物价才开始上涨，一直涨到实际汇率回到原来的水平（本书的推导）。名义汇率与实际汇率在短期里一起动，这是模型可以拿数据检验的地方（\secref{sec:floating-test}）。
% 据: db_eq8

预期在这里是一致的。多恩布什说，完全预见的路径是理性预期在确定情形下的对应物\cite{dornbusch1976}。理性预期在\chref{ch:ratex}。
% 据: db_rational

\begin{example}{货币增加 10\%}\label{ex:floating-overshoot}
设 $\lambda=2$，$\delta=0.5$，$\sigma=1$，$\pi=0.8$，时间以年计（教学用的假设）。（1）按式~\eqref{eq:floating-theta}，$\theta$ 是多少？（2）货币一次增加 10\%，汇率当时上升多少？利率降多少？（3）一年以后，汇率、物价与利率离长期水平各有多远？
\solution
（1）$\sigma/\lambda+\delta=1$，$\theta=0.4+\sqrt{0.16+0.2}=0.4+0.6=1.0$。（2）$\mathrm{d}e/\mathrm{d}m=1+1/(2\times1)=1.5$，汇率上升 15\%，比长期的 10\% 多 5 个百分点。利率降 $10/2=5$ 个百分点，人们预期本币每年升值 $1\times5=5\%$，正好抵消利差。（3）$\mathrm{e}^{-1}\approx0.368$，汇率还高出长期水平 $5\times0.368\approx1.84\%$，物价还低 $10\times0.368\approx3.68\%$，利率还低 1.84 个百分点（本书的计算）。离长期水平的距离每 $\ln2/\theta\approx0.69$ 年减半。
\end{example}

\fig{floating_paths}{货币在第 0 年一次增加 10\% 以后汇率、物价与利率的路径（\exref{ex:floating-overshoot}，本书的计算）。(a) 汇率与物价的对数，以原来的水平为 0；汇率先跳到 15\%，再回落到 10\%，物价从 0 慢慢升到 10\%。(b) 本国利率与外国利率之差（百分点）。参数$\lambda=2$，$\delta=0.5$，$\sigma=1$，$\pi=0.8$是教学用的假设。}{fig:floating_paths}

\figref{fig:floating_paths}画出例题的路径。第 0 年，汇率从 0 跳到 15\%，物价不动，利率差跳到 $-5$ 个百分点。此后汇率回落，物价上涨，利率回升，三者都在两三年里接近新的长期水平（本书的计算）。这期间汇率与物价朝相反的方向走，本币一边升值，物价一边上涨（本书的解读）。

\begin{pitfall}[汇率超调，说明外汇市场不理性]
不一定。在这个模型里，预期与模型自己的路径一致，没有人犯系统的错误，汇率照样超调。原因在商品市场。物价一时不动，货币增加压低了利率，汇率必须先跌过头，人们才会预期本币以后升值，抵消利差（本书的解读）。多恩布什说，他对式（16）的解释还没有用到预期理性的限制\cite{dornbusch1976}。产出若在短期里随需求增加，汇率的贬值可能小于货币的增加，不再超调\cite{dornbusch1976}。
% 据: db_notyet, db_partv
\end{pitfall}

\begin{compare}{固定汇率下的调整与浮动汇率下的调整}{固定汇率}{浮动汇率}
\cmprow{先动的是什么}{储备与货币存量，央行在平价上买卖外汇（本书的概括）}{汇率，央行不必买卖外汇（本书的概括）}
\cmprow{动得多快}{储备在几天里就能大变，日本 11 天买入近 40 亿美元\cite{boj1986}}{汇率当天跳到位，物价要几年才跟上（本书的推导）}
\cmprow{谁先承担}{按旧汇率买入美元的央行与它的国内货币\cite{bis1972,boj1986}}{本币升值时的出口部门\cite{bis1972,epa1972}}
\cmprow{本国的货币政策}{受平价约束，即\chref{ch:mundell}的三选二}{可以自主，瑞士 1973 年浮动以后才完全控制货币供给\cite{rich2003}}
\cmpsum{固定汇率把冲击压在储备与货币上，浮动汇率把它压在汇率上。1971—1973 年各国一再买入美元，最后还是让汇率去动（本书的概括）。}
% 据: boj_11days, b72_exporters, epa_cuts, rich_control
\end{compare}

\section{数据检验与反例}\label{sec:floating-test}

\figref{fig:floating_rates}画出日元、马克、瑞士法郎与英镑对美元的月度汇率。1971 年 8 月以前，四条线大体是平的。日元的月平均一直在 357 到 360 日元之间，英镑、马克与瑞士法郎只有平价调整留下的几次台阶\cite{bisdata46}。1973 年 3 月以后，它们天天在变。日元从 1971 年 7 月的 357.4 日元升到 1973 年 3 月的 261.9 日元，1985 年 2 月又回到 260.4 日元\cite{bisdata46}。
% 据: bisdata_jpy, bisdata_jpy7303, bisdata_jpy8502, bisdata_title

\fig{floating_rates}{日元、马克、瑞士法郎与英镑对美元的汇率，1965—1985 年，月平均，对数刻度，以 1970 年平均为 100（每美元合多少本币，向下表示本币升值）。三条竖线是 1971 年 8 月 15 日、1971 年 12 月 18 日与 1973 年 3 月 19 日。马克按欧元序列乘以 1.95583 换算。数据：国际清算银行数据门户，双边汇率（WS\_XRU）。}{fig:floating_rates}

汇率比物价动得快多少？本书按\chref{ch:external}的实际汇率，以美国为外国，算了月度变化的标准差。名义汇率、实际汇率与两国物价之比都取对数，再算每月的变化\cite{bisdata46}。比较的是 1960—1970 年与 1974—1985 年，中间的过渡年份不算（本书的计算）。每月的对数变化乘以 100，约等于百分比变化，标准差越大，一个月里的涨跌越没有准（本书的解读）。
% 据: bisdata_jpy, bisdata_cpi

日元的变化最典型。1960—1970 年，日元名义汇率每月变化的标准差不到 0.1 个百分点，1974—1985 年是 2.8 个百分点（本书的计算）。实际汇率从 0.8 升到 2.9，两国物价之比前后都是 0.8 左右（本书的计算）。
% 据: bisdata_jpy, bisdata_cpi

马克、瑞士法郎与英镑也一样（\figref{fig:floating_volatility}）。马克名义汇率的标准差从 0.7 升到 2.7。瑞士法郎从 0.2 升到 3.2（本书的计算）。英镑从 0.9 升到 2.6，它在 1960 年代的数字大，是因为 1967 年 11 月贬了值（本书的解读）。去掉 1967 年 11 月与 12 月这两个月，英镑只有 0.1（本书的计算）。
% 据: bisdata_dem, bisdata_chf, bisdata_gbp

浮动以后，四种货币实际汇率的波动是此前的 2.5 到 6.5 倍，名义汇率与实际汇率的波动几乎一样大，物价之比的波动却变化不大（本书的计算）。名义汇率一动，实际汇率就跟着动，物价跟不上，这正是超调模型的前提（本书的解读）。物价之比的波动为什么前后差不多？两国的通胀率在 1970 年代都升高了，可通胀是一年年慢慢变的，月度的变化仍然平稳（本书的解读）。汇率却可以一个月变好几个百分点。
% 据: bisdata_jpy, bisdata_dem, bisdata_chf, bisdata_gbp

\fig{floating_volatility}{名义汇率、实际汇率与两国物价之比（对数）月度变化的标准差（百分点），1960—1970 年与 1974—1985 年（本书的计算）。实际汇率按\chref{ch:external}的式~\eqref{eq:external-real}，以美国为外国。数据：国际清算银行数据门户，双边汇率（WS\_XRU）与长期消费者价格（WS\_LONG\_CPI）。}{fig:floating_volatility}

浮动以后的起落很大。日元 1978 年 10 月升到 183.8 日元，是这十几年里的最高点，1982 年 10 月又跌到 270.9 日元\cite{bisdata46}。按日元计，这四年里美元涨了约 47\%（本书的计算）。瑞士法郎 1978 年 10 月升到 1.54 法郎，正是瑞士国家银行改盯汇率的时候\cite{bisdata46}。英镑 1985 年 2 月跌到每英镑约合 1.10 美元（本书的计算）。美元 1980 年代前期的上涨在\chref{ch:credibility}。
% 据: bisdata_jpy, bisdata_chf, bisdata_gbp

固定汇率年代的名义波动也不全是零。马克 1961 年与 1969 年升值，每次调整都在月度数据里留下一个大跳动（本书的解读）。去掉 1961 年 3 月、1969 年 10 月与 11 月，马克只有 0.2（本书的计算）。浮动以后的波动却是月月都有。1960—1970 年的 132 个月里，四种货币月度变动超过 1\% 的最多只有 3 个月（本书的计算）。1974—1985 年的 144 个月里，这样的月份有 99 到 109 个（本书的计算）。
% 据: bisdata_gbp, bisdata_dem

能不能预测？Meese 与 Rogoff 1983 年比较了各种汇率模型的样本外预测\cite{meeserogoff1983}。他们用 1973 年 3 月到 1981 年 6 月的月度数据，从 1976 年 11 月起逐月预测美元对英镑、马克、日元与贸易加权的汇率\cite{meeserogoff1983}。结论是，在 1 到 12 个月的期限上，随机游走“与任何估计的模型一样好”\cite{meeserogoff1983}。样本外，是说用预测时手头已有的数据估计模型，再拿后来的实际汇率检验，不让模型偷看答案（本书的解读）。随机游走就是拿今天的汇率预测将来的汇率。
% 据: mr_sample, mr_first, mr_asgood, mr_rw

参赛的模型不少。结构模型有 Frenkel—Bilson 的弹性价格货币模型、多恩布什—Frankel 的黏性价格模型，以及加上经常账户的 Hooper—Morton 模型\cite{meeserogoff1983}。时间序列模型与向量自回归也没有更好，远期汇率的预测也不比随机游走强\cite{meeserogoff1983}。
% 据: mr_models, mr_var, mr_forward

他们的检验还对模型格外宽容。结构模型要用到将来的货币、产出与利率，他们给的是实际实现的数值\cite{meeserogoff1983}。均方根误差是预测误差平方的平均再开方，越小越好（本书的解读）。\tabref{tab:floating-mr}列出他们表 1 里 12 个月的预测误差，三个结构模型都比随机游走大\cite{meeserogoff1983}。表里唯一胜过随机游走的，是马克 12 个月的远期汇率\cite{meeserogoff1983}。随机游走的误差以美元对日元最大，一个月是 3.7\%，12 个月是 18.3\%\cite{meeserogoff1983}。
% 据: mr_realized, mr_t1, mr_fwdmark, mr_rw

\begin{table}[tbp]\centering\small
\caption{Meese 与 Rogoff 的样本外预测：12 个月的均方根误差（\%），1976 年 11 月—1981 年 6 月}\label{tab:floating-mr}
\begin{tabularx}{\linewidth}{@{}lYYYYY@{}}
\toprule
汇率 & 随机游走 & 远期汇率 & Frenkel—Bilson & 多恩布什—Frankel & Hooper—Morton\\
\midrule
美元/马克 & 12.98 & 12.60 & 16.12 & 18.87 & 15.69\\
美元/日元 & 18.31 & 18.95 & 18.55 & 20.41 & 19.20\\
美元/英镑 & 9.96 & 11.62 & 14.62 & 13.66 & 14.57\\
\bottomrule
\end{tabularx}
\datasource{Meese 与 Rogoff（1983）第 13 页的表 1 中 12 个月的数字，按页面图像读出。后三列是结构模型，解释变量用的是实现值。表 1 的单变量自回归与向量自回归两列没有列出，它们的误差也都比随机游走大。}
% 据: mr_t1, mr_t1_title, mr_t1_dem, mr_t1_jpy, mr_t1_gbp, mr_purged
\end{table}

两个结果放在一起看。超调模型说明汇率为什么动得比物价大，却没有说明它往哪边动（本书的解读）。\chref{ch:external}讲过非抛补利率平价的失败（\secref{sec:external-test}），也是同一个难处。汇率像资产价格，由对将来的预期决定，预期一变就跳（本书的解读）。
% 据: mr_asgood

\section{争论}\label{sec:floating-debate}

\subsubsection{史密森汇率为什么只维持了十四个月}

一种说法看各国的调整。日本银行在行史里说，布雷顿森林体系缺少顺畅的调整机制，赤字国家以美国为首没有维持国际收支的纪律\cite{boj1986}。日本等盈余国家又拒绝弹性地升值\cite{boj1986}。BIS 1973 年批评美国在 1971 年 8 月 15 日以后“不现实的拒绝”，不肯让美元贬值\cite{bis1973}。尼克松在史密森学会说，原来的汇率因为不稳定而注定失败\cite{nixon1971b}。新汇率是不是又一次调得不够，正是争论所在（本书的解读）。
% 据: boj_defect, boj_discipline, boj_surplus, b73i_unrealistic, nxs_doomed

另一种说法看投机。日本银行认为，史密森体制有两个新弱点，一是美国没有承诺恢复兑换，二是战后从未改过的黄金价格提高了\cite{boj1986}。这给了各国投机者强烈的动机，成了史密森体制的“致命伤”\cite{boj1986}。它还说，1972 年一开年，投机者就发起了攻击\cite{boj1986}。日本政府在史密森协定以后定下绝对避免日元再次升值的方针，到 1973 年 2 月，当局在外汇市场上买入的美元共达 75 亿美元\cite{boj1986}。BIS 则说，汇率调整一再拖延，总要等到市场危机才启动，史密森以后也是一样\cite{bis1973}。
% 据: boj_weak, boj_fatal, boj_attack, boj_avoid, boj_75, b73_delayed

美国的方案本身也有分寸问题。美国要的改善是 130 亿美元，国际组织算的是 90 亿美元，罗马会议以后美国同意考虑更少的数（\secref{sec:floating-smithsonian}）\cite{bis1972}。调得少，是不是注定撑不久，还是撑不久另有原因，史料没有一致的说法。至今没有定论。
% 据: b72_13, b72_9, b72_rome

\subsubsection{超调能不能解释浮动汇率的波动}

多恩布什承认，市场调整速度的差别缺少理论依据，但它推出的汇率行为与最近的经验相似\cite{dornbusch1976}。本书算出的实际汇率与名义汇率一起放大，也合乎物价调整慢的前提（\secref{sec:floating-test}）。
% 据: db_backing, bisdata_jpy

多恩布什说，商品市场调整慢，并没有很有说服力的理论支持，但事实明显指向这个方向\cite{dornbusch1976}。在他的模型里，调整越慢，贸易条件的变化越大，也越持久\cite{dornbusch1976}。
% 据: db_facts, db_terms

Meese 与 Rogoff 的结果却是，即使给了实现的解释变量，结构模型也没有胜过随机游走\cite{meeserogoff1983}。他们说，结构模型的预测已经排除了解释变量的不确定性，表现不佳多半要归于联立方程偏误、抽样误差、真实参数的随机变动或者设定错误\cite{meeserogoff1983}。他们列了这些可能的原因，没有得出确定的结论\cite{meeserogoff1983}。石油价格冲击与宏观政策体制的改变带来的结构不稳定，模型没有充分纳入的其他实际扰动，都可能重要\cite{meeserogoff1983}。至今没有定论。
% 据: mr_realized, mr_noconclusion, mr_oil

\subsubsection{浮动该不该有人管}

\secref{sec:floating-smithsonian}讲过 BIS 1972 年列举的理由，各国不愿让本币完全自由浮动\cite{bis1972}。1973 年 3 月 16 日，十国集团同意必要时由央行干预来管理浮动\cite{bis1973}。美国一方的考虑却在政治。3 月 14 日，基辛格在电话里对财政部副部长西蒙（William E. Simon，1927—2000）说，他的理由“完全出于政治”\cite{frus46b,loc46}。
% 据: b72_reasons, b73_g10, f33_date, f33_political, loc_simon

基辛格说，要尽量阻止欧洲形成统一的立场\cite{frus46b}。他说，若是支持某一种单独的货币，他会更愿意\cite{frus46b}。15 日晚上，他又说，应当让共同浮动“尽可能难以运转”，两人同意不干预合乎这个目标\cite{frus46b}。西蒙补充说，或者有时干预，帮一些人，不帮另一些人\cite{frus46b}。可是不管也有代价。BIS 1976 年说，1975—1976 年央行的干预达到了让人想起盯住汇率年代的规模\cite{bis1976}。
% 据: f33_prevent, f33_individual, f35_hard, f35_nonint, f35_some, b76_intervention

瑞士的经历站在两边之间。瑞士国家银行 1978 年回头盯住过汇率（\chref{ch:monetarism}）。1990 年代中期，它记着 1978—1981 年的难处，一再说维护物价稳定的任务与公众要求压低法郎实际汇率的呼声可能冲突\cite{rich2003}。Rich 事后说，它若承认自己定货币政策时考虑了汇率，会更准确地描述实际的做法\cite{rich2003}。至今没有定论。
% 据: rich_80, rich_1995, rich_conflict, rich_admit

\begin{deeper}[史料说明]
尼克松 1971 年 8 月 15 日的讲话与 12 月 18 日对十国集团部长的讲话，据美国总统项目（American Presidency Project）网站转载的《总统公文集》。FRUS 是美国国务院历史学家办公室编的档案集，本章读了 1969—1976 年第 3 卷的第 164、165、168、211、221 号与第 31 卷的第 12、14、33、35 号文件，都是网页版。录音与会议笔记都按编者的转述，基辛格的回忆录也只按编者的转引。

BIS 的年报读的是 1972、1973、1974 与 1976 年的扫描本与识别文字层。识别文字认错的地方照录，并在快照的说明里一一列出，如把 2¼ 认成“2\%”，把 4½ 认成“414”。1973 年年报的几页文字层字母间夹着空格，第 136 页的 4½ 与 2¼ 又认错了，这些都按页面图像读出，\tabref{tab:floating-smithsonian}的数字也按页面图像核对过。日本银行的《日本银行百年史》第 6 卷没有文字层，按渲染的页面图像逐字读出。

几个日子与幅度在史料里不一致。第二次贬值与共同浮动的日子，正文已经说明，本书按各自的说法写出，不作取舍。共同浮动的窄幅，BIS 1973 年的年报说是名义交叉汇率上下各 2.25\%，日本银行的行史说是上下各 1.125\%，正文用的是 BIS 的数字。佐佐木直的生卒年在美国国会图书馆的名称规范档里没有找到，那里的 Sasaki, Tadashi 是另一个人（佐々木正）。

多恩布什与 Meese—Rogoff 的论文仍在版权期内，读的是网上的第三方副本，只用来核对，不随本书分发。多恩布什的公式按页面图像读出，文字层把希腊字母认错了。Meese 与 Rogoff 1981 年的讨论稿是同一研究的早期版本，封面要求引用前先征得作者同意，本书只把它列为早期版本\cite{meese1981}，结果都取自 1983 年发表的版本。表 1 的数字按页面图像读出。

欧洲共同体理事会 1972 年 3 月 21 日的决议读的是《欧洲共同体公报》的法文本，中文是本书的翻译。汇率与消费者价格取自 BIS 数据门户。联邦德国的汇率在那里按欧元折算，本书乘以 1.95583 还原成马克，还原后的数字与\tabref{tab:floating-smithsonian}马克的中心汇率及幅度相合（本书的核对）。1960 年代各国的消费者价格指数常常几个月不变，英国 1960—1970 年的 132 个月度变化里有 76 个是零，物价之比的月度变化因此时有时无（本书的核对）。
% 据: nxs_note, nx_note, f168_diary, b72i_margins, b73i_rule, b73_margin, boj_march11, ifdp_cover, ifdp_land, mr_journal, db_journal, oj_date, bisdata_dem, bisdata_cpi
\end{deeper}

\section{今天还剩什么}\label{sec:floating-legacy}

牙买加的第四条留了下来，各国至今自选汇率安排，有的浮动，有的盯住，有的在两者之间（本书的概括）。蛇形浮动可以看作欧洲货币体系的前身（本书的解读），后来走到了欧元，在\chref{ch:euro}。主要货币之间的大起大落也没有停，1985 年的广场协议在\chref{ch:japanbubble}。
% 据: b76_article4, b73i_snake

1971 年那种一个星期就买进几十亿美元的事，在大国之间少了，央行的干预却一直都在（本书的概括）。BIS 1976 年看到的就是这样一种局面，规则正式结束了，干预的规模却像盯住汇率的年代\cite{bis1976}。浮动与盯住之间，各国一直在找位置。
% 据: b76_end, b76_intervention

浮动也给了各国自己管货币的余地，这正是\chref{ch:mundell}的三选二。Rich 说，瑞士改为浮动以后，央行在原则上完全控制了货币供给，并用它来对付 1973 年初已升到约 7\% 的通胀\cite{rich2003}。\chref{ch:monetarism}讲的货币目标，就是在这个前提下定的（本书的解读）。
% 据: rich_control, rich_7

浮动以后，汇率的风险从央行转到了企业与投资者身上。1971 年日本的出口产地靠压价与加班消化升值，浮动以后，这样的调整天天都要做（本书的解读）。
% 据: epa_cuts

超调模型成了讲汇率的标准一课，Meese 与 Rogoff 的随机游走成了检验汇率模型的标尺（本书的解读）。两者并不矛盾。模型说明，物价调整慢，货币冲击就会被汇率放大；它没有说冲击什么时候来，往哪边来（本书的解读）。
% 据: db_abstract, mr_asgood

\begin{takeaways}
\begin{itemize}
\item 1971 年 8 月 15 日，美国停止美元兑换黄金，加征 10\% 的进口附加税。决定在戴维营作出，国务院与国家安全委员会没有参加。
\item 史密森协定让美元对黄金贬值 7.9\%，波动幅度放宽到平价上下各 2.25\%。它只维持了十四个月，1973 年 3 月主要货币转向浮动，1976 年的牙买加协议承认了这个局面。
\item 多恩布什的超调模型说，物价调整慢，货币增加时汇率先跳过长期水平，$\mathrm{d}e/\mathrm{d}m=1+1/\lambda\theta$，然后随物价上涨而回落。
\item 浮动以后，名义汇率与实际汇率的月度波动放大了几倍到几十倍，物价之比的波动变化不大。Meese 与 Rogoff 发现，汇率模型的预测不比随机游走好。
\item 欧洲用蛇形浮动维持成员国之间的窄幅，日本、瑞士与联邦德国的央行在浮动前后都大量买入过美元。
\end{itemize}
\end{takeaways}

\begin{selfcheck}
\begin{enumerate}
\item 按\exref{ex:floating-overshoot}的参数，若 $\lambda=4$ 而 $\theta=1$，货币增加 8\%，汇率当时上升多少？利率降多少？
\item 按\exref{ex:floating-overshoot}的结果，两年以后汇率还高出长期水平多少？
\item 史密森协定以后，日元对美元的中心汇率是 308 日元，比原来的 360 日元升值多少？
\item 美元对黄金贬值 7.9\%，黄金的美元价格上升多少？
\item 1971 年 12 月到 1973 年 2 月是几个月？
\end{enumerate}
\end{selfcheck}

\begin{answers}
\item $1+1/(4\times1)=1.25$，汇率上升 $8\times1.25=10\%$，利率降 $8/4=2$ 个百分点。
\item $5\times\mathrm{e}^{-2}\approx0.68\%$。
\item $360/308-1\approx16.9\%$，与表中对美元的升值相同。
\item $1/(1-0.079)-1\approx8.6\%$，即从每盎司 35 美元升到 38 美元。
\item 14 个月。
\end{answers}

回到 1971 年 8 月 15 日的晚上。在美国买美国货的人，手里的美元确实没有因此少买东西。尼克松在 1973 年 2 月还提醒部下强调这一点\cite{frus46b}。到 1973 年 3 月，同样一美元在东京只能换 260 多日元，而不是 360 日元\cite{bisdata46}。日本银行在行史里说，把储备大多放在美元上的产油国等，很可能卖出美元\cite{boj1986}。产油国若以美元收款，美元贬值就等于石油卖得便宜了（本书的解读）。油价冲击为什么同时推高通胀和失业？产油国为什么会被“资源诅咒”？
% 据: nx_tomorrow, f12_note, bisdata_jpy, boj_360, boj_oil
